% TOP_PROBLEM: {"problemId": "problem.freyd-s-generating-hypothesis", "canonicalTitle": "Freyd's Generating Hypothesis", "releaseRank": 482, "primaryDomain": "geometry_topology", "primaryDomainLabel": "Geometry & topology", "displayStatus": "open", "releaseStatus": "open"}
% !TeX program = lualatex
% Definition and sources only; this file is not an LLM solution attempt.
\documentclass[11pt]{article}
\usepackage[margin=25mm]{geometry}
\usepackage{fontspec}
\setmainfont{Latin Modern Roman}
\usepackage{amsmath,amssymb,mathtools}
\usepackage{unicode-math}
\setmathfont{Latin Modern Math}
\usepackage{xurl,hyperref}
\hypersetup{colorlinks=true,urlcolor=blue}
\setcounter{secnumdepth}{0}
\setlength{\parindent}{0pt}
\setlength{\parskip}{5pt}
\setlength{\emergencystretch}{3em}
\begin{document}
\section*{\texorpdfstring{Freyd's Generating Hypothesis}{Freyd's Generating Hypothesis}}\label{482}
\subsection{Definitions and mathematical statement}
A finite spectrum is a spectrum built from finitely many cells, or equivalently an object of the finite stable homotopy category, with the usual retract convention. Write $[X,Y]$ for stable homotopy classes of maps and $\pi_nX=[\Sigma^nS,X]$. The graded group $\pi_*X$ is a module over the stable homotopy ring $\pi_*S$.

Freyd's generating hypothesis says that for finite spectra $X,Y$,
\[
 \bigl(\pi_n(f)=0\text{ for every }n\in{\mathbb{Z}}\bigr)
                  \Longrightarrow f=0\text{ in }[X,Y].
\]
Equivalently, $[X,Y]\to\operatorname{Hom}_{\pi_*S}(\pi_*X,\pi_*Y)$ is injective, with degree-zero graded homomorphisms on the right. The target is faithfulness, not surjectivity of this map. Allowing arbitrary infinite spectra changes the statement.
\par\smallskip\textit{Formulation and concept references:} \href{https://math.uchicago.edu/~may/PEOPLE/AMB/Hovey_OnFreydsGenHyp.pdf}{[S1]}.

\subsection{Short English statement}
Is a map between finite spectra necessarily null-homotopic when it acts as zero on every stable homotopy group?
\subsection{Sources}
\begin{itemize}
\item[S1]Mark Hovey, On Freyd's Generating Hypothesis, Quarterly Journal of Mathematics 58 (2007), 31–45.
\url{https://math.uchicago.edu/~may/PEOPLE/AMB/Hovey_OnFreydsGenHyp.pdf}.
\textit{Formulation source.}
\end{itemize}
\end{document}
